\section{Three-row generating functions}\label{app:threegf}

We finish the proof of Theorem~\ref{thm:threegf}, with $D=LN=1-2z-4z^3-4z^4$. A
\emph{restricted run} is a string of singletons with no two letters $b$ at distance two.
Reversal and reflection carry the next lemma to $BA$, $BC$ and $CB$.

\begin{lemma}[forbidden windows in the classes containing $B$]\label{lem:threewindows}
In the class $AB$ and the triple class, a window fails exactly when the table lists it for its
class, with $s$ any singleton.
\begin{center}\small
\begin{tabular}{@{}lll@{}}
\toprule
class & windows & forbidden \\
\midrule
$AB$, $X$ & singletons only & $bsb$ \\
$AB$ & one double & $bsA$, $Asb$, $bAb$, $bAc$, $cAb$, $aBa$, $aBb$, $bBa$, $cBc$ \\
 & two doubles & $bAB$, $ABa$ \\
$X$ & one triple & $bsX$, $Xsb$, $aXa$, $aXb$, $bXa$, $bXc$, $cXb$, $cXc$ \\
\bottomrule
\end{tabular}
\end{center}
\end{lemma}

\begin{proof}
Singleton windows follow from Lemma~\ref{lem:dichotomy}, and $sAt$ as in
Lemma~\ref{lem:sixofnine}(a), except that $bAb$ now fails test (e) of Lemma~\ref{lem:diamond}
as $B\in F$. In $sBt$ and $sXt$ both corners in column $i+1$ are black, so the
diamond is empty for $s=t=b$, and for outer $s$ and $t$ test (c) fails it exactly when $s=t$.
For $s=b\ne t$ it shows $(i+2,2)$ alone, and test (e) fails it exactly when $\{2,t\}\in F$. In
$uAB$ test (d) fails $u=b$, as $(i+2,3)$ and $(i+1,2)$ are black, and otherwise three corners
show. In $ABw$ the diamond is empty for $w=b$, and test (e) fails it exactly for $w=a$. In $Asb$
and $Xsb$ tests (b) and (e) fail the vertical pair or the point of column $i+1$. Other windows
with a double or $X$ at the left tip show three corners or an oblique pair needing the white
point $(i+1,2)$. The rest is symmetric.
\end{proof}

Give an outer letter bit $0$ and weight $2$, and give $b$ bit $1$. Let $h_{ij}$ count, by
length, the continuations of a restricted run ending in bits $i,j$, where $b$ cannot follow
$1,j$. Read runs $w_1\cdots w_n$ with $n\ge1$
outward from a double. Let $g_0$ and $g_1$ count those with $w_1$ a prescribed outer letter and
with $w_1=b$, and $u$ and $v$ the same runs with $w_2$, if present, outer. A fictitious letter
before $w_1$, outer for $g$ and $b$ for $u$ and $v$, gives $(g_0,g_1,u,v)=z(h_{00},h_{01},
h_{10},h_{11})$. Solving the linear system for $h$ gives
\begin{equation}\label{eq:runs}
Dg_0=z(1+z+3z^2+2z^3),\ \ Dg_1=z(1+z),\ \ Du=z(1+2z^2+2z^3),\ \ Dv=z.
\end{equation}

\begin{proposition}[the triple class]\label{prop:GX}
With $u$ and $v$ as in~\eqref{eq:runs}, $G_X(z)=z\bigl(1+4u+2v+2u^2+v^2\bigr)$, and its
denominator divides $D^2$.
\end{proposition}

\begin{proof}
In $UXW$ the runs are restricted, and $bsX$ and $Xsb$ make the second letters from $X$ outer.
One nonempty run gives $2u+v$ on either side. With two, the window at $X$ allows only the
neighbor pairs $(a,c)$, $(b,b)$ and $(c,a)$, which give $2u^2+v^2$.
\end{proof}

A word of $AB$ is $UAVBW$ with restricted runs $U$, $V$ and $W$. By
Lemma~\ref{lem:threewindows}, a nonempty $U$ read from $A$ is counted by $u$ or $v$, and a
nonempty $W$ read from $B$ by $g_0$ or $g_1$. Let $\ell_t$ and $r_t$ count the allowed $U$ and
$W$, empty included, when $t$ follows $A$ or precedes $B$. The windows at the
double give
\begin{equation}\label{eq:lr}
(\ell_a,\ell_b,\ell_c)=(1+2u+v,1+u,1+2u),\ (r_a,r_b,r_c)=(1+g_0,1+g_0+g_1,1+g_0+g_1).
\end{equation}
For $s\in\{00,10\}$ let $J_{s\beta}$ count the continuations from $s$ ending in bit $\beta$, the
empty one with $\beta=0$. Solving as for $h$ gives, with rows $s$ and columns $\beta$,
\begin{equation}\label{eq:J}
J=\frac1D\begin{pmatrix}1+2z^2+2z^3&z+z^2\\1&2z^2+2z^3\end{pmatrix}.
\end{equation}

\begin{proposition}[the class $\{A,B\}$, polynomial form computer-assisted]\label{prop:GAB}
With~\eqref{eq:runs} to~\eqref{eq:J},
\begin{equation}\label{eq:GAB}
\begin{aligned}
G_{AB}(z)={}&z^2(1+2u)(1+g_0+g_1)+z^3(\ell_ar_a+\ell_br_b+\ell_cr_c)\\
&+z^4\,(\ell_a+\ell_c,\ \ell_b)\,J\,(r_a+r_c,\ 2r_b)^{T}.
\end{aligned}
\end{equation}
Moreover $G_{AB}=P_B/D^3$ with
\begin{equation}\label{eq:PB}
P_B(z)=z^2(1+z)(1+z^2-3z^3+4z^6+4z^7).
\end{equation}
\end{proposition}

\begin{proof}
Split by the length $k$ of $V$. For $k=0$ the windows $uAB$ and $ABw$ give the first term, as
$u\ne b$ and $w\ne a$. For $k=1$ the window $AtB$ passes, and $uAt$ and $tBw$ give $\ell_tr_t$.
For $k\ge2$ the window $Av_1v_2$ makes $v_2$ outer, $v_{k-1}v_kB$ passes, and $v_3\cdots v_k$
is a continuation from bits $00$ or $10$, as $v_1$ is outer or $b$, with weight
$\ell_{v_1}r_{v_k}$. Since $r_a+r_c$ counts both outer $v_k$, for $k\ge3$ the weight $2$ of an outer $v_k$
in $J$, like the $2$ in $2r_b$, counts the choices of $v_2$, and $k=2$ is the constant term of
$J$.
Exact computer algebra gives~\eqref{eq:PB}.
\end{proof}

\begin{theorem}[the exact denominator, computer-assisted]\label{thm:denominator}
The reduced denominator of $\sum_{m\ge3}a_3(m)z^m$ is $Q=R^3L^3N^3$.
\end{theorem}

\begin{proof}
By Theorem~\ref{thm:threegf}(a) the denominator divides $R^3D^3=Q$. No two of $R$, $L$ and $N$
share a root, as $D(\frac13)\ne0$ and $L-N=-2z-4z^2$ vanishes only where $N\ne0$. At
$z=\frac13$ only $G_{AC}$ has a pole, and the numerator in~\eqref{eq:GAC} is $\frac4{81}$. As $P_B\bmod L=\frac z2$ and $P_B\bmod N=\frac14-\frac z2$, at each root
of $D$ the term $G_{AB}$ has a triple pole, $G_X$ at most a double pole, and $G_{AC}$ none.
\end{proof}
